%% A280_Landauer_Phasenraum_Buchhaltung_De.tex
%% FFGFT A-Serie -- A280: Landauer und die Phasenraum-Buchhaltung
%% Autor: Johann Pascher  --  ORCID 0009-0000-6518-4064
%% Sprache: Deutsch  --  Engine: LuaLaTeX
%% Quelle: Arbeitsnotiz "Landauer_Phasenraum_Buchhaltung_final.md", 24. Juli 2026

\documentclass[11pt,a4paper]{article}

% ---- Engine: LuaLaTeX ----
\usepackage{fontspec}
\usepackage{mathtools}
\usepackage{unicode-math}
\usepackage{polyglossia}
\setmainlanguage{german}

\usepackage{microtype}
\emergencystretch=3.5em
\tolerance=2500
\hbadness=3000

\usepackage{enumitem}
\usepackage{booktabs}
\usepackage{longtable}
\usepackage{array}
\usepackage[a4paper,margin=2.6cm,headsep=1.1em,headheight=22pt]{geometry}
\usepackage{fancyhdr}
\usepackage{titlesec}
\usepackage{xcolor}
\usepackage{needspace}
\usepackage{tocloft}

% ---- Fonts ----
\setmainfont{Inter}[
  Scale=1.02,
  Ligatures=TeX ]
\setsansfont{Inter}[Scale=MatchLowercase,Ligatures=TeX]
\setmonofont{JetBrains Mono}[Scale=0.88]
\setmathfont{Libertinus Math}

% ---- Colours ----
\definecolor{accent}{HTML}{1F3A5F}     % deep slate blue
\definecolor{accent2}{HTML}{6A4E2E}    % muted bronze
\definecolor{rulecol}{HTML}{B9C2CE}
\definecolor{linknav}{HTML}{2C5F8A}
\definecolor{linkcite}{HTML}{7A5C2E}

% ---- Schichtmarkierungen (A-Serie) ----
\definecolor{laySet}{HTML}{7A2E2E}
\definecolor{layK}{HTML}{1D5C3A}
\definecolor{layB}{HTML}{1F3A5F}
\definecolor{layS}{HTML}{6A4E2E}
\definecolor{layH}{HTML}{5A2E6A}

\newcommand{\LS}{\textcolor{laySet}{\textbf{\footnotesize\textsf{[SETZUNG]}}}\,}
\newcommand{\LK}{\textcolor{layK}{\textbf{\footnotesize\textsf{[K]}}}\,}
\newcommand{\LB}{\textcolor{layB}{\textbf{\footnotesize\textsf{[B]}}}\,}
\newcommand{\LSk}{\textcolor{layS}{\textbf{\footnotesize\textsf{[S]}}}\,}
\newcommand{\LH}{\textcolor{layH}{\textbf{\footnotesize\textsf{[H]}}}\,}

\newcommand{\pruef}[1]{{\small\itshape(Prüfskript, #1)}}

% ---- Hyperref ----
\usepackage[
  colorlinks=true,
  linkcolor=linknav,
  citecolor=linkcite,
  urlcolor=linknav,
  bookmarksnumbered=true,
  bookmarksopen=true,
  pdfstartview=FitH ]{hyperref}
\hypersetup{
  pdftitle={Landauer und die Phasenraum-Buchhaltung (A280)},
  pdfauthor={Johann Pascher},
  pdfsubject={Statistische Mechanik --- Landauer-Grenze, Phasenraum, Gebietsreduktion},
  pdfkeywords={Landauer, Liouville, Phasenraum, Entropie, Gebietsreduktion, FFGFT, A280}
}

% ---- Section styling ----
\titleformat{\section}
  {\color{accent}\normalfont\Large\bfseries}
  {\color{accent2}\thesection}{0.7em}{}
\titleformat{\subsection}
  {\color{accent}\normalfont\large\bfseries}
  {\color{accent2}\thesubsection}{0.6em}{}
\titlespacing*{\section}{0pt}{1.6\baselineskip}{0.7\baselineskip}
\titlespacing*{\subsection}{0pt}{1.1\baselineskip}{0.4\baselineskip}

% ---- Headers ----
\pagestyle{fancy}
\fancyhf{}
\renewcommand{\headrulewidth}{0.4pt}
\renewcommand{\headrule}{\hbox to\headwidth{\color{rulecol}\leaders\hrule height \headrulewidth\hfill}}
\fancyhead[L]{\footnotesize\color{accent}\itshape FFGFT $\cdot$ A280 $\cdot$ Landauer und die Phasenraum-Buchhaltung}
\fancyhead[R]{\footnotesize\color{accent2}\thepage}

% ---- TOC look ----
\renewcommand{\cftsecfont}{\normalfont}
\renewcommand{\cftsecpagefont}{\normalfont}
\renewcommand{\cftsecleader}{\cftdotfill{\cftdotsep}}
\setlength{\cftbeforesecskip}{3pt}

\setlength{\parskip}{0.35em}
\setlength{\parindent}{0pt}
\linespread{1.06}

\begin{document}

% ===================== TITEL =====================
\begin{titlepage}
\centering
\vspace*{2.2cm}
{\color{rulecol}\rule{0.5\textwidth}{0.6pt}}\\[1.4cm]
{\color{accent}\Huge\bfseries Landauer und die\\[0.3em] Phasenraum-Buchhaltung\par}
\vspace{1.0cm}
{\large\itshape Warum Gebietsreduktion prinzipiell Wärme kostet\\ --- und warum der Inhalt egal ist\par}
\vspace{0.6cm}
{\color{accent2}\rule{0.28\textwidth}{0.4pt}}\\[1.0cm]
{\large FFGFT A-Serie $\cdot$ A280\par}
\vspace{0.5cm}
{\normalsize Standalone-Fassung mit Schichtmarkierungen, Prüfskript-Bezug\\ und vollständiger Quellenliste\par}
\vfill
{\large Johann Pascher\par}
\vspace{0.3cm}
{\footnotesize ORCID: 0009-0000-6518-4064\par}
\vspace{0.2cm}
{\footnotesize Juli 2026\par}
\vspace{1.4cm}
{\color{rulecol}\rule{0.5\textwidth}{0.6pt}}
\vspace*{1.5cm}
\end{titlepage}

% ===================== TOC =====================
\tableofcontents
\thispagestyle{fancy}
\clearpage

% ===================== KURZFASSUNG =====================
\begingroup\small\noindent
\textbf{Kurzfassung.} Die Landauer-Grenze wird hier nicht als Aussage über
\emph{Information} entwickelt, sondern als Aussage über \emph{Gebiete im
Phasenraum}. Ein physikalisches System belegt ein Gebiet; eine Abbildung, die
mehrere Ausgangsgebiete auf ein Zielgebiet reduziert, verkleinert dieses
Volumen; der Satz von Liouville verbietet Volumenschrumpfung im abgeschlossenen
System; also muss ein Bad das Volumen übernehmen, und der Preis ist Wärme
$Q \ge k_B T \ln(W_{\text{v}}/W_{\text{n}})$ --- mit $\ln 2$ nur als
Binär-Spezialfall. Gezählt wird nach Unterscheidbarkeit (Orthogonalität), nicht
nach Energie; die absolute Grund-Diskretisierung liefert $h^f$. Daraus folgt
unmittelbar die Inhaltsneutralität: die Thermodynamik zählt Gebiete und liest
sie nicht. Drei Ebenen werden strikt getrennt --- Konstruktion (Ingenieur),
Zustandsübergang (Takt), Beschreibung (Programm) ---, und daraus folgt, dass ein
realer Rechner nicht diskret Bits löscht, sondern einen kontinuierlichen
Entropiestrom trägt (DRAM-Refresh, Qubit-Dekohärenz). Der epistemische Vorbehalt
wird ausdrücklich gebucht: nicht messbar heißt nicht nicht vorhanden --- Landauer
gilt wegen der Mechanik (Liouville), nicht wegen Unwissenheit; die
Irreversibilität ist statistisch, nicht ontologisch absolut. \emph{Berührung mit
dem $\xi$-Schema der FFGFT: keine} --- die Landauer-Buchhaltung lebt auf der
Ebene der statistischen Mechanik (eigene Ebene, eigene Buchhaltung).
\par\endgroup

\bigskip

\begingroup\small\noindent
\textbf{Schichtmarkierungen.} \LS deklarierter Ausgangspunkt ---
\LB algebraisch/mathematisch aus deklarierten Grundlagen abgeleitet ---
\LK empirisch geprüft bzw.\ aus Primärquellen belegt ---
\LSk Plausibilitätsskizze ---
\LH offene Forschungsfrage (Erweiterung des A-Serien-Schemas, siehe
Anhang~\ref{sec:korr}, Punkt~6).
\par\endgroup

\bigskip

\begingroup\small\noindent
\textbf{Prüfskript.} \texttt{landauer\_pruefskript.py} (Checks 1--10). Alle
numerischen Aussagen dieses Dokuments sind dort mit Assertions hinterlegt.
Check~7: Zustandszahl im Kolloid-Gebiet; Check~8: Analogfall mit deklariertem
$\delta x$.
\par\endgroup

\bigskip

% ===================== 1 =====================
\section{Das Gebiet kommt zuerst}

Zwei Gläser stehen auf einem Tisch --- Gebiet A und Gebiet B. Daneben das
Weltmeer. Man schüttet eines der Gläser in das Meer. Das Glasvolumen verschwindet
nicht --- es verteilt sich auf 1,4 Milliarden Kubikkilometer. Dieses Bild trägt
das ganze Dokument: Es gibt keine verschwindende Information, keine vernichteten
Bits --- nur eine Gebietsreduktion, die einen messbaren Preis hat.

\LS \textbf{Ein physikalisches System belegt ein Gebiet im Phasenraum.} Dieses
Gebiet hat eine Größe --- messbar in Einheiten von $h^f$ (Abschnitt~\ref{sec:zaehlbasis}).
Es existiert unabhängig davon, wie man es benennt, einteilt oder beschreibt. Das
ist der physikalische Ausgangspunkt. Alles andere kommt danach.

\LB \textbf{Beschreibungen kommen von außen.} Wenn man entscheidet, das Gebiet in
zwei Teilgebiete einzuteilen und diese „A“ und „B“ zu nennen und das Zielgebiet
„Z“, hat man eine Beschreibung der Gebietsreduktion eingeführt. Nennt man sie
stattdessen „0“, „1“ und „0“ (traditionelle Notation), entsteht die Verwechslung
mit logischen Bits --- diese Konvention verdeckt, dass „0“ sowohl Ausgangsgebiet
als auch Zielgebiet bezeichnet, und suggeriert einen Informationsinhalt, wo keiner
ist. Wenn man feiner einteilt --- in 1000 Zonen ---, hat man ${\sim}10$ logische
Bit eingeführt. Wenn man nur eine Zahl zwischen 0 und 1 zuordnet, hat man einen
analogen Wert eingeführt. Die Physik erzwingt keine dieser Entscheidungen. Das
Gebiet ist in allen Fällen dasselbe physikalische Gebiet.

Die Begriffe „Bit“, „Byte“, „Qubit“, „analoger Wert“ sind Bezeichnungen für
Beschreibungsentscheidungen --- nicht für physikalische Objekte. Sie teilen das
Gebiet ein, sie erschaffen es nicht.

\subsection{Das Landauer-Bit ist Hardware --- nicht Information}

\LB Das wird regelmäßig verwechselt. Klarstellung in einem Absatz: Ein
Landauer-Bit ist ein physikalisches Gebiet mit zwei stabilen Makrozuständen ---
ein Transistor, ein Kondensator, eine Leiterbahn, die hoch oder niedrig liegt. Ein
Bus-Bit ist dasselbe: ein physikalisches Gebiet mit zwei stabilen Makrozuständen,
eingebettet in eine parallele Hardware-Architektur. Eine Speicherzelle ist
dasselbe. Die drei sind strukturell identisch. Die Thermodynamik behandelt alle
drei gleich:
\[
  Q \;\ge\; k_B\,T\,\ln\!\left(\frac{W_{A+B}}{W_Z}\right)
\]
--- unabhängig davon, ob das Gebiet „Landauer-Bit“, „Bus-Bit“ oder
„Speicherzelle“ heißt.

\textbf{Die Hardware ist festgelegt, bevor gerechnet wird. Sie ändert sich nicht
durch das, was man mit ihr rechnet.} Ein 64-Bit-Bus hat 64 physikalische Gebiete
--- vor der Rechnung, während der Rechnung, nach der Rechnung. Was die Rechnung
mit diesen Gebieten macht --- welche Abbildungen sie ausführt, welche Gebiete auf
welche Zielgebiete reduziert werden --- ist eine logische Beschreibung dieser
physikalischen Vorgänge. Die Beschreibung existiert nicht in der Hardware; sie ist
eine externe Interpretation.

Der Informationsgehalt eines Bits --- ob es „wahr“ oder „falsch“ bedeutet, ob es
zu einem korrekten Rechenschritt gehört, ob es überhaupt gelesen wird --- ist eine
Entscheidung auf der Beschreibungsebene. Die Thermodynamik kennt diese Ebene
nicht. Sie kennt nur: Gebiet A und Gebiet B existieren; nach der Gebietsreduktion
existiert Zielgebiet Z; $W_Z \le W_{A+B}$; das Bad nimmt den Unterschied auf.

\subsection{Drei Ebenen --- strikt getrennt}

\LB \textbf{Ebene 1 --- Konstruktion.} Der Ingenieur entscheidet beim Chip-Design,
welche Gebietsoperationen existieren --- bijektiv oder nicht-bijektiv. Das ist in
Silizium gegossen, unveränderlich während des Betriebs, vor jeder Rechnung. Diese
Entscheidung legt die thermodynamische Kostenstruktur fest --- nicht die Rechnung,
nicht das Programm.

\LB \textbf{Ebene 2 --- Zustandsübergang im laufenden Betrieb.} Sobald der Rechner
läuft, finden Zustandsübergänge taktweise statt --- gleichmäßig, mit der
Taktfrequenz, unabhängig davon, was berechnet wird. Landauer-Kosten fallen dabei
nur bei \textbf{nicht-bijektiven} Gebietsabbildungen an --- wenn ein Bauteil durch
seinen physikalischen Aufbau mehrere Eingangsgebiete auf ein Zielgebiet Z
reduziert. Das ist durch die Hardware-Architektur festgelegt, nicht durch eine
Rechenoperation. Ein Transistor, der von Gebiet A nach Gebiet B kippt, führt eine
bijektive Abbildung aus --- kein Volumenverlust, keine Landauer-Kosten. Was die
Beschreibungsebene danach mit dem Zustand des Bauteils macht, ist thermodynamisch
irrelevant.

\textbf{Die Landauer-Kosten laufen mit der Taktfrequenz, praktisch konstant,
solange der Rechner läuft} --- weil die nicht-bijektiven Gebietsabbildungen durch
den Konstrukteur festgelegt sind und mit jedem Taktzyklus anfallen. Die
CPU-Auslastung sagt wenig darüber aus: Sie misst, wie viele Taktzyklen für
Berechnungen auf der Beschreibungsebene genutzt werden --- aber im Leerlauf laufen
die Taktübergänge weiterhin durch. Der thermodynamische Unterschied zwischen hoher
und niedriger Auslastung ist gering, solange keine physikalischen Teile
tatsächlich abgeschaltet werden. Clock Gating, Power Gating --- das sind
Hardware-Entscheidungen auf Ebene~1, keine Programm-Entscheidungen. Die Hardware
bestimmt --- nicht das Programm.

\LB \textbf{Ebene 3 --- Beschreibung.} Das Programm, die Rechnung, eine
KI-Steuerung können entscheiden, welche Hardware genutzt wird --- Teile
abschalten, Routen wählen, bestimmen, dass für ein bestimmtes Ergebnis die volle
Bitbreite nicht benötigt wird. Das hat nichts mit Landauer zu tun. Es ist
Ressourcenmanagement auf der Beschreibungsebene. Die thermodynamischen Kosten der
genutzten Hardware fallen trotzdem an --- festgelegt durch Konstruktion und
Schaltmoment, unabhängig davon, was die Beschreibungsebene mit den Ergebnissen
macht.

\subsection{Wie viel steckt in einem physikalischen Gebiet?}

\LK Ein Kolloid-Teilchen in einer 1-\textmu m-Falle enthält physikalisch rund
$\mathbf{7{,}9\cdot10^{9}}$ \textbf{orthogonale Quantenzustände}
\pruef{Check 7}. Nennt man grob zwei Zonen „links“ und „rechts“, beschreibt man
dieses Gebiet mit 1 logischem Bit. Teilt man es in 1000 Zonen, beschreibt man
dasselbe Gebiet mit ${\sim}10$ logischen Bit. Das physikalische Gebiet ist in
beiden Fällen identisch --- die Beschreibungstiefe ist eine Wahl des Beobachters.

\LB \textbf{Die Kosten des Löschens hängen am Gebiet, nicht an der Beschreibung.}
Die Formel lautet
\[
  \boxed{\;Q \;\ge\; k_B\,T\,\ln\!\left(\frac{W_{\text{vorher}}}{W_{\text{nachher}}}\right)\;}
\]
$W_{\text{vorher}}$ und $W_{\text{nachher}}$ sind physikalische Gebietsgrößen in
Einheiten von $h^f$. Wie man dieses Verhältnis benennt --- ob als „1 Bit löschen“
oder „Faktor 2 Reduktion“ --- ändert nichts an der Physik.

\LK \textbf{Der analoge Fall zeigt das besonders klar --- und zugleich seine
Komplexität.} Wie viele unterscheidbare Zustände ein analoges Signal enthält,
hängt von mehreren Faktoren ab, die alle physikalisch sind: von der Bandbreite
$B$, der Temperatur $T$ und dem Widerstand $R$ (Johnson-Nyquist-Rauschen,
$U_{\text{rms}} = \sqrt{4 k_B T R B}$ [B20]), von der Empfindlichkeit der
Messtechnik, vom Signal-Rausch-Abstand des konkreten Systems und von weiteren
externen Störquellen. Es gibt keine universelle Auflösungsgrenze für analoge
Signale --- sie ist eine Systemeigenschaft, keine Naturkonstante.

Was sich aber allgemein sagen lässt: Wenn die Auflösungsgrenze durch das
thermische Rauschen des Bades selbst gesetzt wird (der physikalisch tiefste Fall),
dann gilt das bereits in Abschnitt~\ref{sec:zaehlbasis} formulierte Prinzip:
Dasselbe Bad, das die Löschkosten kassiert, begrenzt die Stufung des analogen
Signals. Die Gebietsanzahl $W = L/\delta x$ und die Löschkosten $k_B T \ln W$
skalieren mit derselben Temperaturgröße --- die Buchhaltung ist selbstkonsistent.
Wie viele logische Bit man dem analogen Gebiet zuschreibt, ist davon unabhängig
und bleibt eine Beschreibungsentscheidung \pruef{Check 8: Illustrationsbeispiel
mit deklariertem $\delta x$}.

\subsection{Shannon und Boltzmann sind verschiedene Buchführungen}

\LB [B18][B19] Shannon beschreibt, wie viel Unsicherheit eine \emph{Beschreibung}
trägt: $H = -\sum_i p_i \log_2 p_i$ --- eine Größe über Beschreibungen. Boltzmann
beschreibt, wie groß ein \emph{physikalisches Gebiet} ist: $S = k_B \ln W$ ---
eine Größe über Gebiete. Die beiden sind nur dann identisch, wenn man die
Beschreibung exakt auf die physikalischen Mikrozustände abbildet und diese
gleichverteilt sind. In jedem anderen Fall --- grobe Einteilung, ungleiche
Besetzung, kontinuierliche Systeme --- sind sie verschiedene Dinge mit
verschiedenen Einheiten.

\LK \textbf{Was Landauer 1961 wirklich sagte} [B1]: Er schrieb ausdrücklich
„logically irreversible“ --- logisch irreversibel, nicht physikalisch
irreversibel. Er meinte: eine Abbildung, bei der verschiedene logische
Eingangszustände auf denselben logischen Ausgangszustand führen. Das ist eine
Aussage über Beschreibungen.

\LB \textbf{Und deshalb ist sie für sich genommen keine physikalische Bedingung.}
Hinreichend für die Wärmeproduktion ist allein die \textbf{Nicht-Bijektivität der
Gebietsabbildung}: dass mehrere Ausgangsgebiete auf ein kleineres Zielgebiet
fallen. „Operation“ ist an dieser Stelle das falsche Wort --- eine Operation ist
eine Beschreibung, und Beschreibungen kosten nichts. Mit der Gebietsabbildung
fällt die logische Irreversibilität nur unter einer zusätzlichen, ausdrücklich zu
deklarierenden Realisierungsannahme zusammen: dass die logischen Marken auf
disjunkte physikalische Gebiete abgebildet sind und beim Übergang kein Gebiet
mitgeführt wird. Ist diese Annahme verletzt, trägt die logische Irreversibilität
nichts --- Bennett [B3] führt jede logisch irreversible Funktion als Folge
bijektiver Gebietsabbildungen aus, ohne Untergrenze
(Abschnitt~\ref{sec:gebietsop}). Und umgekehrt kostet jede nicht-bijektive
Gebietsabbildung, ob sie überhaupt als Operation beschrieben wird oder nicht.
Jeder physikalisch irreversible Prozess produziert Wärme, unabhängig davon, ob er
als „logisch irreversibel“ beschrieben wird; das ist genau der Punkt von
Abschnitt~\ref{sec:immer}. Die Physik unter der Beschreibung kann sehr viel
reicher sein als die Beschreibung selbst.

\LB \textbf{Der Fehler der Standarddarstellung} ist damit: Sie behandelt die
Beschreibungseinheit --- das Bit --- so, als wäre sie eine physikalische
Mindestgröße. Das ist sie nicht. Eine physikalische Mindestgröße gibt es: $h^f$.
Aber sie hat nichts mit Bits zu tun. Das Gebiet kommt zuerst; wie viele Bits man
ihm zuschreibt, ist eine nachträgliche Entscheidung.

% ===================== 2 =====================
\section{Die eine Idee, aus der alles folgt}

Ein Bit ist physikalisch: zwei \textbf{unterscheidbare} Zustände eines Systems.
Zwei Mulden für ein Teilchen, zwei Magnetisierungsrichtungen, zwei Spannungspegel.
Entscheidend ist nur: Das System kann in Zustand „0“ oder „1“ sein, und man kann
beide auseinanderhalten.

\textbf{Gebietsreduktion: Egal, in welchem Gebiet --- A oder B --- das System
vorher war, nachher ist es im Zielgebiet Z.} Das Zielgebiet Z ist frei wählbar;
was zählt, ist ausschließlich das Verhältnis $W_{A+B}/W_Z$.

\begin{center}
\begin{minipage}{0.42\textwidth}
\ttfamily\small
\begin{tabbing}
\ \ \ \ A \ ---+\\
\ \ \ \ \ \ \ \ \ \ +--->\ \ Z\\
\ \ \ \ B \ ---+
\end{tabbing}
\end{minipage}
\end{center}

Zwei Gebiete werden auf eines reduziert --- das Zielgebiet Z. Diese Abbildung ist
\textbf{nicht umkehrbar}: Wer nur das Zielgebiet kennt, kann nicht rekonstruieren,
in welchem Ausgangsgebiet das System war. Genau das --- und nichts anderes ---
erzeugt die Kosten (Abschnitt~\ref{sec:liouville}).

% ===================== 3 =====================
\section{Was „Phasenraum“ hier bedeutet}

Der Phasenraum ist der Raum aller möglichen Mikrozustände (alle Orte und Impulse
aller Teilchen). Ein Makrozustand wie „Gebiet A“ ist darin kein Punkt, sondern ein
\textbf{Gebiet}: viele Mikrozustände (thermisches Zittern in der linken Mulde ---
jede Zitterlage zählt als „0“).

\LS Die Boltzmann-Entropie misst dieses Gebiet:
\[
  S \;=\; k_B \ln W
\]
$W$ = Phasenraumvolumen des Gebiets. Entropie ist hier ein \textbf{Maß für die
Größe des Aufenthaltsgebiets} --- nichts Mystisches.

\begin{center}
\begin{minipage}{0.75\textwidth}
\ttfamily\small
\begin{tabbing}
Phasenraum vorher:\ \ \ \ \ \ \ \=Phasenraum nachher:\\[0.3em]
+-------+-------+\>+-------+-------+\\
|\ \ \ \ \ \ \ |\ \ \ \ \ \ \ |\>|\ \ \ \ \ \ \ |\ \ \ \ \ \ \ |\\
|\ \ A\ \ \ \ |\ \ B\ \ \ \ |\ \ -->\ \>|\ \ Z\ \ \ \ | leer\ \ |\\
|\ \ \ \ \ \ \ |\ \ \ \ \ \ \ |\>|\ \ \ \ \ \ \ |\ \ \ \ \ \ \ |\\
+-------+-------+\>+-------+-------+\\
\ Volumen: 2W0 (A+B)\>\ Volumen: W0 (nur Z)
\end{tabbing}
\end{minipage}
\end{center}

\LB Löschen halbiert das zugängliche Volumen:
\[
  \Delta S_{\text{Gebiet}} \;=\; -\,k_B \ln 2
\]
Reine Buchhaltung über Gebietsgrößen --- kein Modell, keine Hardware.
(Verallgemeinerung auf $N$ Stufen und Nicht-Gleichverteilung:
Abschnitt~\ref{sec:zaehlbasis}.)

% ===================== 4 =====================
\section{Warum das nicht gratis geht: Liouville}
\label{sec:liouville}

Die mikroskopische Dynamik (klassisch Hamiltonsch, quantenmechanisch unitär) hat
eine harte Eigenschaft, den \textbf{Satz von Liouville} [B2]:

\begin{quote}
\LB Die Zeitentwicklung eines abgeschlossenen Systems \textbf{erhält das
Phasenraumvolumen}. Verformen, verschieben, verwirbeln --- ja. Schrumpfen --- nie.
\end{quote}

Präzisierung: Liouville gilt exakt nur für das abgeschlossene System \emph{ohne}
Bad. Für das Teilsystem „Bit“ allein --- nach Kopplung ans Bad --- gilt es nicht
mehr. Abgeschlossen ist nur das Gesamtsystem Bit${}+{}$Bad; für dieses bleibt das
Gesamtvolumen erhalten, während das Bit-Teilgebiet sich auf Kosten des Bad-Gebiets
verkleinern darf.

Anschaulich: Der Phasenraumfluss ist eine \textbf{inkompressible Flüssigkeit}. Ein
Tintenklecks im Wasser lässt sich beliebig verzerren; sein Volumen bleibt.

Damit schnappt die Falle zu:
\begin{itemize}[leftmargin=1.4em,itemsep=0.2em]
  \item Löschen \textbf{verlangt}: $2W_0 \to W_0$ (Halbierung, Abschnitt~3).
  \item Liouville \textbf{verbietet}: Volumenschrumpfung im abgeschlossenen System.
\end{itemize}

\LB \textbf{Folgerung: Ein abgeschlossenes System kann nicht löschen.} Kein
Ingenieursproblem --- ein Satz. Der einzige Ausweg: Das System darf nicht
abgeschlossen sein. Das Volumen muss irgendwohin --- ins Bad
(Abschnitt~\ref{sec:bad}).

% ===================== 5 =====================
\section{Das Bad als Volumen-Empfänger}
\label{sec:bad}

Koppelt man das Bit an ein Wärmebad, gilt Liouville für das \textbf{Gesamtsystem}
Bit${}+{}$Bad. Soll das Bit-Gebiet auf die Hälfte schrumpfen
(Gleichverteilungsfall: beide Gebiete gleich groß), muss das Bad-Gebiet sich
mindestens \textbf{verdoppeln}. Im allgemeinen Fall gilt
$\Delta S_{\text{Bad}} \ge -\Delta S_{\text{Bit}}$ (zweiter Hauptsatz); die
Verallgemeinerung auf $\ln(W_{\text{vorher}}/W_{\text{nachher}})$ folgt in
Abschnitt~\ref{sec:zaehlbasis}:
\[
  \Delta S_{\text{Bad}} \;\ge\; +\,k_B \ln 2
\]

Ein Bad bei Temperatur $T$ vergrößert sein Gebiet durch Energieaufnahme; die
Grundbeziehung $dS = \delta Q/T$ übersetzt das in Wärme:
\[
  \LB\quad Q \;\ge\; T\,\Delta S_{\text{Bad}} \;=\; k_B\,T \ln 2
\]

Das ist die Landauer-Grenze [B1][B3]:

\begin{quote}
\textbf{Liouville (Volumen bleibt) $+$ Buchhaltung (Bit-Gebiet halbiert)
$\Rightarrow$ Bad-Gebiet verdoppelt $\Rightarrow$ Wärme $\ge k_B T \ln 2$ ins
Bad.}
\end{quote}

Bei 300\,K: $k_B T \ln 2 \approx 2{,}87\cdot10^{-21}$\,J pro Bit
\pruef{Check 1}.

Direkt ablesbar die zwei Bedingungen:
\begin{enumerate}[leftmargin=1.6em,itemsep=0.2em]
  \item \textbf{Ohne Bad keine Löschung} --- es gibt keinen anderen
        Volumen-Empfänger (Abschnitt~\ref{sec:liouville}).
  \item \textbf{Gleichheit nur quasistatisch} --- jede endliche Geschwindigkeit
        erzeugt Zusatzdissipation $W \approx k_B T \ln 2 + B/\tau$ [B4][B5]
        (Abschnitt~\ref{sec:experimente}).
\end{enumerate}

% ===================== 6 =====================
\section{Warum eigentlich 2? --- Die Zählbasis}
\label{sec:zaehlbasis}

Die Binarität ist \textbf{Konvention, nicht Physik}. Allgemein:
\[
  Q \;\ge\; k_B\,T\,\ln\!\left(\frac{W_{\text{vorher}}}{W_{\text{nachher}}}\right)
\]
--- ein Verhältnis von Gebietsgrößen. $N$ Stufen auf eine gelöscht: $\ln N$
\pruef{Check 4}. Die $\ln 2$ ist der Informatik-Spezialfall.

\LB \textbf{Was zählt als Stufe? Nicht die Energie.} Energieniveaus taugen nicht
als Zählbasis: Sie können entartet sein oder ungleiche Sprünge haben, ohne dass
sich die Zählung ändert. Das Kriterium ist \textbf{Unterscheidbarkeit}: Zwei
Zustände zählen getrennt, wenn sie im Prinzip fehlerfrei auseinanderzuhalten sind.
Für \textbf{Quantensysteme} heißt das Orthogonalität im Hilbertraum; für
\textbf{klassische Systeme} heißt es Unterscheidbarkeit innerhalb der durch das
Bad gesetzten thermischen Auflösungsgrenze --- d.\,h.\ der Abstand der Zustände
muss größer sein als die thermischen Fluktuationen. Ein Bit kann in zwei
entarteten Zuständen gleicher Energie liegen --- ideal sogar: Halten kostet dann
nichts, nur Vergessen. Niveau-Abstände wirken nur indirekt, über die thermischen
Besetzungswahrscheinlichkeiten. Die eigentliche Währung ist daher die Gibbs-Form
[B6]
\[
  S \;=\; -\,k_B \sum_i p_i \ln p_i
\]
mit $\ln N$ als Spezialfall der Gleichverteilung. Ungleich gestufte, ungleich
besetzte Systeme tragen automatisch weniger Entropie \pruef{Checks 5--6}.

\LB \textbf{Die Grund-Diskretisierung existiert --- und ist $h^f$.} Klassisch hat
das Phasenraumvolumen Einheiten; $W$ ist nur bis auf eine willkürliche Zellgröße
definiert. Die Buchhaltung stört das nicht, weil nur Differenzen eingehen --- die
Zelle kürzt sich. Die Quantenmechanik setzt die Zelle absolut: Die Zahl
orthogonaler Zustände in einem Gebiet $\Gamma$ eines Systems mit $f$
Freiheitsgraden ist
\[
  N_{\max} \;=\; \Gamma / h^f \qquad \text{[B7]}
\]
\pruef{Check 7: Kolloid in 1-\textmu m-Falle, rund $7{,}9\cdot10^{9}$ orthogonale
Zustände --- die „2“ des Experiments ist grob gewählt, nicht fundamental}.

\subsection{Fundamentale vs.\ emergente Zweiwertigkeit}

\begin{center}
\small
\begin{tabular}{@{}p{0.20\textwidth}p{0.30\textwidth}p{0.38\textwidth}@{}}
\toprule
 & \textbf{Spin-½} & \textbf{Kolloid in Doppelmulde} \\
\midrule
Herkunft der 2 & fundamental: Hilbertraum $\mathbb{C}^2$ & emergent: 2 metastabile Makro-Gebiete \\
Gültigkeit & immer & nur solange Barriere $\gg k_B T$ auf der Prozesszeitskala \\
darunter & nichts --- keine feinere Stufung & Quasikontinuum von Mikrozuständen \\
\bottomrule
\end{tabular}
\end{center}

\subsection{Der quasi-analoge Fall}

Ein Wert im Bereich $L$ mit Auflösung $\delta x$ trägt $\ln(L/\delta x)$ Stufen;
Gebietsreduktion kostet $k_B T \ln(L/\delta x)$ \pruef{Check 8}. $\delta x$ ist
nicht frei: Das Badrauschen setzt es --- dasselbe Bad, das die Löschkosten
kassiert, begrenzt die Stufung. Kein Analog-Trick führt an der Grenze vorbei.

% ===================== 7 =====================
\section{Warum der Inhalt egal ist}
\label{sec:inhalt}

Die gesamte Rechnung hat an \textbf{keiner Stelle} gefragt, welche Bedeutung
Gebiet A oder B trägt, ob das Bit in einer korrekten Rechnung steht oder nicht, ob
es als „wahr“ oder „falsch“ interpretiert wird. Eingegangen ist ausschließlich:
\textbf{wie viele} unterscheidbare Gebiete vorher, \textbf{wie viele} nachher.

Die Thermodynamik zählt Gebiete; sie liest sie nicht. Gebiet A und Gebiet B sind
physikalisch gleichwertig --- die Physik kennt keine bevorzugte Richtung, und das
Zielgebiet Z ist frei wählbar. Welche Bedeutung man ihnen gibt („wahr“/„falsch“,
0/1, hoch/niedrig), ist eine Entscheidung außerhalb der Thermodynamik:

\begin{quote}
\LB \textbf{Gebietsreduktion kostet für jede Belegung identisch --- unabhängig
davon, wie man die Gebiete interpretiert.} Es gibt keinen thermodynamischen
Inhalts- oder Wahrheitsdetektor.
\end{quote}

In derselben Sprache: Eine nicht-terminierende Berechnung dissipiert unbegrenzt
Energie und hat dabei keinen Wahrheitswert. Energieumsatz und semantische
Gültigkeit liegen auf verschiedenen Ebenen --- Phasenraum-Buchhaltung hier,
Repräsentation dort. Kategorientrennung: physikalische Gültigkeit
(Hardware-Ebene) vs.\ repräsentationale Gültigkeit (Inhaltsebene).

\subsection{Hardware und Rechenoperation sind verschiedene Dinge}

Ein grundlegender Denkfehler in der Standarddarstellung von Landauer: Die Hardware
wird als Abbild der Rechenoperation behandelt --- als ob der physikalische Apparat
nur dann existiert und nur dann thermodynamisch wirksam ist, wenn gerade eine
logische Operation ausgeführt wird. Das stimmt nicht.

\LB \textbf{Die Hardware ist nicht die Rechenoperation.} Ein 64-Bit-Prozessor ist
ein physikalisches System mit 64-Bit-Zustandsraum. Dieser Zustandsraum existiert
unabhängig davon, ob gerade eine sinnvolle Operation läuft, ob das Ergebnis 1 Bit
oder 64 Bit braucht, welche Bedeutung das Ergebnis trägt, ob jemand das Ergebnis
liest oder nicht. Die Hardware ist keine Abstraktion --- sie ist Silizium,
Leitungen, Felder, Kondensatoren, die kontinuierlich mit ihrer Umgebung
wechselwirken.

\LB \textbf{Die Rechenoperation ist eine Abstraktion über der Hardware.} Sie
beschreibt, was der Programmierer oder Mathematiker sieht --- Eingabe, Ausgabe,
Funktion. Die Hardware kennt diese Abstraktion nicht. Sie kennt nur
Zustandsübergänge, Ladungsverschiebungen, Wärmeabgaben.

\textbf{Die Konsequenz für Landauer.} Wenn man fragt „was kostet das Löschen eines
Bits“, nimmt man die Abstraktion als Ausgangspunkt und sucht die thermodynamischen
Kosten der logischen Operation. Das ist eine legitime Frage --- aber sie verfehlt
die Physik des realen Systems. Der reale Rechner produziert Entropie
kontinuierlich, unabhängig von der logischen Abstraktion, die gerade auf ihm
läuft.

\subsection{Dieselbe Aussage von der Hardware-Seite}

Die Aussage aus Abschnitt~\ref{sec:inhalt} --- die Thermodynamik liest keinen
Inhalt --- lässt sich unabhängig von der Buchhaltung direkt aus der physikalischen
Struktur des Rechners herleiten.

\LB \textbf{Der Rechner schrumpft nicht.} Ein 64-Bit-Prozessor hat vor und nach
jeder Operation gleich viele Transistoren, gleich viele Leitungen, gleich viele
Freiheitsgrade. Der Phasenraum der Hardware bleibt konstant --- $2^{64}$ mögliche
Zustände vorher, $2^{64}$ nachher. Das Ergebnis braucht vielleicht 1 Bit; der
Apparat, der es produziert hat, ist trotzdem gleich groß geblieben.

\LB \textbf{Was wirklich ins Bad geht, ist die Korrelation.} Die logische
Abbildung ist viele-zu-eins: viele verschiedene Eingangszustände landen auf
demselben Ausgang. Die 63 weggeworfenen Bits sind logisch weg, aber physikalisch
steckt ihre Spur noch irgendwo in den Transistorzuständen --- nur nicht mehr aus
dem Ergebnis rekonstruierbar. Genau diese Nicht-Rekonstruierbarkeit, die
Korrelation zwischen Eingang und Ausgang, die im Ergebnis nicht mehr steckt, geht
als Wärme ans Gitter. Nicht das Silizium, sondern das Vergessen der Korrelation.

\LB \textbf{Der Vergessensvorgang ist inhaltsneutral.} Ein Prozessor, der
$2+2=4$ berechnet, und einer, der $2+2=5$ berechnet, schalten dieselbe Anzahl
Transistoren, werfen dieselbe Bitbreite weg und zahlen dieselbe thermodynamische
Rechnung ans Bad. Die Physik unterscheidet die beiden nicht. Ein Ergebnis in
Gebiet A kostet dasselbe wie eines in Gebiet B --- die Physik unterscheidet die
Gebiete nicht.

Das ist keine neue Behauptung --- es ist eine Herleitung von unten. Landauers
Abstraktion „Gebietsreduktion kostet, Inhalt egal“ folgt aus der konkreten Physik
des Rechners, ohne Formel. Die Buchhaltung in Abschnitt~\ref{sec:inhalt} und die
Hardware-Struktur hier kommen auf verschiedenen Wegen zur selben Aussage.

\subsection{Der Rechner rechnet immer}
\label{sec:immer}

Ein physikalischer Rechner hört nicht auf zu rechnen, wenn er auf ein Ergebnis
wartet. Die Taktfrequenz läuft, die Transistoren schalten, das Netzteil liefert
Strom --- kontinuierlich, unabhängig davon, ob gerade „nützliche“ Arbeit
passiert. Ein Idle-Prozessor rechnet Leerschleifen. Ein wartender Server schickt
Keepalive-Pakete. Ein Betriebssystem pollt Interrupts. Es gibt keinen Zustand
„Rechner rechnet nicht“, solange er eingeschaltet ist.

Das bedeutet: Der Vergessensvorgang läuft \textbf{immer}. Die thermodynamische
Rechnung --- Korrelationsverlust, Wärme ans Gitter --- hört nicht auf, egal ob das
Ergebnis gebraucht wird oder nicht, egal welche Bedeutung es trägt.

Das verschärft Landauers Formulierung zur Irrelevanz in der Praxis. Er isoliert
einen einzelnen diskreten Gebietsreduktions-Vorgang --- „Gebietsreduktion kostet
$k_B T \ln 2$“. Die Realität ist ein \textbf{kontinuierlicher Strom} von
Zustandsübergängen, Korrelationsverlusten und Wärmeabgaben, drei Milliarden Mal
pro Sekunde bei einem modernen Prozessor, ohne Pause.

\textbf{Konkretes Beispiel: Dynamisches RAM.} Ein DRAM-Bit ist ein Kondensator,
der Ladung hält. Aber Kondensatoren verlieren Ladung durch Leckströme ---
kontinuierlich, unaufhaltsam. Deshalb muss jede DRAM-Zelle alle 32--64
Millisekunden \textbf{aufgefrischt} werden: Der Refresh-Vorgang liest den Zustand
und schreibt ihn zurück. Das passiert taktweise, automatisch, für jede der
Milliarden Zellen --- egal ob die Daten gerade gebraucht werden, egal welche
Bedeutung sie tragen, egal ob das Programm gerade schläft.

Das ist kein Nebeneffekt --- das ist die physikalische Natur des Systems. Der
DRAM-Refresh ist \textbf{kein} Gebietsreduktions-Vorgang im Landauer-Sinne; er ist
ein Erhaltungsvorgang, und seine direkten Kosten sind primär Ingenieurskosten
(Leckströme, Umladung). Aber der Grund, warum er überhaupt stattfinden muss, ist
thermodynamisch: Der Kondensator koppelt permanent ans Wärmebad. Die Ladung ---
das physikalische Gebiet „Bit geladen“ --- wird durch die thermische Kopplung
kontinuierlich in Richtung des gemeinsamen Gleichgewichts gedrängt. Das ist eine
permanente, unaufhörliche Wechselwirkung mit dem Bad, unabhängig davon, ob gerade
eine logische Operation stattfindet. Der Refresh ist die technische Antwort auf
diese thermodynamische Wirklichkeit --- er renormiert das Gebiet zurück, weil das
Bad es permanent verwischt.

Aber das ist noch nicht der ganze Punkt. Der Refresh \textbf{selbst} ist keine
passive Reaktion --- er erfordert aktive Rechenarbeit: Ein Zähler muss laufen, die
Adresse der nächsten zu refreshenden Zeile muss berechnet werden, ein
Timing-Signal muss erzeugt werden, der Speicherbus muss gesperrt werden, der
Leseverstärker muss aktiviert werden. Das ist echte Rechenarbeit --- mit echten
Zustandsübergängen, echten Korrelationsverlusten, echtem Entropiestrom ins Bad.
Und der Speicher-Controller gibt dafür in einer Bank mit 8192 Zeilen und 64\,ms
Refresh-Intervall rund $1{,}3\cdot10^{5}$ Refresh-Befehle pro Sekunde aus --- ohne
Unterbrechung, völlig unabhängig davon, ob das laufende Programm gerade etwas
berechnet oder schläft.

Das ist der eigentliche Punkt: Nicht der Refresh-Vorgang selbst ist
Landauer-relevant im Sinne einzelner Bit-Löschungen, sondern die Tatsache, dass
\textbf{der Rechner laufend rechnen muss, um seinen eigenen Zustand zu erhalten}
--- Refresh-Zähler, Takt-Generatoren, Watchdog-Timer,
Betriebssystem-Scheduler, Speicher-Controller. Das ist eine permanente
Eigenberechnung, die nichts mit dem Nutzprogramm zu tun hat. Die Thermodynamik
dieser Eigenberechnung beschreibt Landauer nicht, weil er einzelne Operationen
betrachtete; es ist eine zusätzliche, fundamentalere Kostenkategorie:
\textbf{Der Rechner zahlt Entropie dafür, dass er existiert --- nicht nur dafür,
dass er rechnet.}

\LK Eine DRAM-Zelle, die eine Stunde lang „nichts tut“, hat in dieser Stunde
trotzdem rund $1{,}12\cdot10^{5}$ Refresh-Zyklen durchlaufen (bei 32\,ms
Intervall; bei 64\,ms rund $5{,}6\cdot10^{4}$) und dabei kontinuierlich Entropie
produziert. Die gespeicherte Information war die ganze Zeit dieselbe --- aber die
thermodynamische Rechnung lief durch.

Das ist der physikalische Rechner, wie er wirklich ist: kein Automat, der diskret
Bits löscht, sondern ein System, das kontinuierlich gegen den thermodynamischen
Zerfall ankämpft und dafür bezahlt. Die Frage lautet daher nicht: Wie viel kostet
die Gebietsreduktion? Sondern: Wie groß ist der kontinuierliche Entropiestrom
eines physikalischen Rechners als Funktion seiner Taktfrequenz, seiner Bitbreite
und seiner logischen Irreversibilität --- unabhängig davon, was er gerade
berechnet?

Das ist eine offene Frage. Fredkin, Toffoli und Landauer selbst haben darüber
nachgedacht [B3][B13]. Die Standarddarstellung setzt immer noch am einzelnen
Bit-Gebietsreduktions-Vorgang an, weil das handhabbar ist. Das ist eine
Abstraktion, die die Physik des realen Rechners verfehlt: Der Rechner ist kein
Automat, der auf Befehl einzelne Bits löscht --- er ist ein physikalisches System,
das kontinuierlich Entropie produziert, ob man hinschaut oder nicht, ob das
Ergebnis stimmt oder nicht.

\subsection{Quantencomputer: Dekohärenz als kontinuierlicher Entropiestrom}

Beim Quantencomputer gilt dasselbe Prinzip wie beim DRAM --- der physikalische
Apparat kämpft kontinuierlich gegen den thermodynamischen Zerfall --- aber die
Physik ist unerbittlicher.

\textbf{Was Dekohärenz ist.} Ein Qubit in Superposition hat eine Phasenbeziehung
zu seiner Umgebung. Jede kleinste Wechselwirkung --- ein Phonon, ein Photon, eine
Magnetfeldfluktuation --- verändert diese Phase. Das Qubit verliert seine
Quantennatur nicht, weil Information vernichtet wird, sondern weil die Korrelation
zwischen Qubit und Umgebung ins Bad verteilt wird. Wieder Liouville: Das Volumen
bleibt erhalten, aber es verteilt sich auf Qubit-plus-Umgebung-Korrelationen, die
niemand mehr kontrolliert.

\textbf{Die Verschärfung gegenüber DRAM.} Beim klassischen DRAM kann man
refreshen --- Zustand lesen, zurückschreiben, kostet thermodynamisch,
funktioniert, weil 0 und 1 robust sind. Beim Qubit ist das prinzipiell unmöglich:
Eine Messung kollabiert die Superposition. Einen unbekannten Quantenzustand kann
man weder kopieren (No-Cloning-Theorem [B14]) noch messen, ohne ihn zu verändern.
Refresh im klassischen Sinne gibt es nicht.

Was es stattdessen gibt: \textbf{Quantenfehlerkorrektur} [B15]. Aber die ist
thermodynamisch noch teurer als DRAM-Refresh --- sie braucht 100--1000
physikalische Qubits pro logischem Qubit, kontinuierliche Syndrom-Messungen und
klassische Korrekturoperationen. Das alles läuft durch, taktweise, für jedes
logische Qubit, auch wenn gerade keine Gate-Operation stattfindet.

\LK \textbf{Dekohärenzzeit als DRAM-Analogon.} Ein supraleitendes Qubit hat
typisch $T_2 \sim 100\,\text{\textmu s}$ [B16] --- danach ist die
Phaseninformation im Bad verteilt. Das ist rund 640-mal kürzer als ein
DRAM-Refresh-Intervall von 64\,ms, aber qualitativ anders: Der DRAM verliert
Information langsam und stetig, das Qubit verliert sie exponentiell schnell, und
die „Information“, die verloren geht, ist nicht ein Bit, sondern eine
kontinuierliche komplexe Amplitude über einem hochdimensionalen Hilbertraum (in
der Praxis durch $h^f$ begrenzt, Abschnitt~\ref{sec:zaehlbasis}).

\textbf{Der Entropiestrom beim Quantencomputer.} Beim klassischen Rechner ist der
kontinuierliche Entropiestrom durch Taktfrequenz und Bitbreite bestimmt. Beim
Quantencomputer kommt ein zusätzlicher Term dazu: der dekohärenz-induzierte
Entropiestrom, der proportional zur Anzahl der Qubits und umgekehrt proportional
zu $T_2$ läuft --- unabhängig davon, ob gerade ein Gate ausgeführt wird oder
nicht. Ein wartender Quantencomputer produziert Entropie aggressiver als ein
wartender klassischer Rechner, weil die Quantennatur selbst gegen die Umgebung
ankämpft.

\textbf{Dekohärenz ist keine Gebietsreduktion im Landauer-Sinne} --- sie ist ein
Verlust von Quantenkorrelationen. Aber thermodynamisch zahlt sie denselben Preis:
Die Korrelation geht ins Bad, das Bad nimmt das Volumen auf. Der entscheidende
Unterschied zum DRAM: Beim DRAM kann man die Information zurückgewinnen
(Refresh), beim Qubit nicht (No-Cloning-Theorem [B14]). Dekohärenz und
Landauer-Löschen sind verschiedene Prozesse mit demselben thermodynamischen
Mechanismus.

\textbf{Inhaltsneutral wie immer.} Dekohärenz unterscheidet nicht zwischen einem
Qubit, das einen bestimmten Quantenzustand trägt, und einem, das einen anderen
trägt. Sie zerstört beide gleich schnell. Der thermodynamische Preis des
Quanten-Rechnens hängt nicht am Inhalt --- er hängt an der Physik des Systems.

\LH \textbf{Offene Frage.} Beim klassischen Rechner ist der kontinuierliche
Entropiestrom im Prinzip durch reversible Logik reduzierbar
(Abschnitt~\ref{sec:gebietsop}). Beim Quantencomputer ist das fundamentaler
begrenzt: Ob es einen prinzipiellen Minimalpreis für das Aufrechterhalten von
Quantenkohärenz gibt --- unabhängig von der Güte der Fehlerkorrektur --- ist eine
offene Frage der Quantenthermodynamik [B17].

% ===================== 8 =====================
\section{Gebietsoperationen --- volumenerhaltend oder volumenverkleinernd}
\label{sec:gebietsop}

Jede physikalische Zustandsänderung ist eine \textbf{Abbildung zwischen
Phasenraumgebieten}. Die thermodynamisch entscheidende Frage ist ausschließlich:
Erhält die Abbildung das Gebietsvolumen, oder verkleinert sie es?

Ob diese Abbildung zu einer „Rechnung“ gehört, ob sie als logische Operation
beschrieben wird, ob jemand sie als sinnvoll bezeichnet --- das ist für die
Thermodynamik vollständig irrelevant.

\LB \textbf{Volumenerhaltende Abbildungen} --- das Gebiet wird umgeformt, nicht
verkleinert. Liouville gilt. Keine thermodynamische Untergrenze:

\begin{center}
\begin{minipage}{0.75\textwidth}
\ttfamily\small
A --->\ B\ \ \ \ (bijektiv: jedes Eingangsgebiet auf ein eigenes Ausgangsgebiet)\\
B --->\ A
\end{minipage}
\end{center}

Traditionell heißen solche Abbildungen „reversible Gatter“ (NOT, CNOT, Toffoli)
--- physikalisch sind es \textbf{bijektive Gebietsabbildungen}. Ob sie Teil einer
Rechnung sind oder nicht, ist irrelevant.

\LB \textbf{Volumenverkleinernde Abbildungen} --- mehrere Eingangsgebiete landen
auf demselben Zielgebiet Z. Volumen schrumpft. Das Bad muss den Unterschied
aufnehmen. Thermodynamische Untergrenze:

\begin{center}
\begin{minipage}{0.75\textwidth}
\ttfamily\small
AA ---+\\
AB ---+---> Z\ \ \ \ \ \ 4 Eingangsgebiete -> 2 Ausgangsgebiete\\
BA ---+\ \ \ \ \ \ \ \ \ \ \ \ \ Volumen schrumpft => das Bad zahlt\\
BB -------> Z
\end{minipage}
\end{center}

Traditionell heißen solche Abbildungen „irreversible Gatter“ (AND, OR, ERASE) ---
physikalisch sind es \textbf{nicht-bijektive Gebietsabbildungen}. Wieder: ob sie
Teil einer Rechnung sind oder nicht, ist irrelevant.

\LB \textbf{Bennett} [B3] hat gezeigt, dass jede Folge von Gebietsabbildungen im
Prinzip als reine Folge volumenerhaltender Abbildungen ausgeführt werden kann ---
Zwischenzustände werden mitgeführt und am Ende reversibel freigegeben. Die
thermodynamischen Kosten fallen nicht bei den Zwischenschritten an, sondern
ausschließlich beim endgültigen Verzicht auf Gebiete:

\begin{quote}
\LB \textbf{Die Untergrenze klebt nicht an der Abbildung, sondern am Verzicht auf
Gebiete --- am Vergessen.}
\end{quote}

% ===================== 9 =====================
\section{Die Buchhaltung in einer Tabelle}

\begin{center}
\small
\begin{tabular}{@{}llll@{}}
\toprule
\textbf{Schritt} & \textbf{Bit-Gebiet} & \textbf{Bad-Gebiet} & \textbf{Gesamt (Liouville [B2])} \\
\midrule
vorher & $2W_0$ & $W_{\text{Bad}}$ & $2W_0 \cdot W_{\text{Bad}}$ \\
Reduktion auf Z & $W_Z \le W_0$ & ? & muss $\ge 2W_0 W_{\text{Bad}}$ bleiben \\
also & $W_Z$ & $\ge 2\,W_{\text{Bad}}$ & erhalten (ja) \\
Preis & $\Delta S = -k_B \ln\frac{W_{\text{v}}}{W_Z}$ & $\Delta S \ge +k_B \ln\frac{W_{\text{v}}}{W_Z}$ & $Q \ge k_B T \ln\frac{W_{\text{v}}}{W_Z}$ ins Bad \\
\bottomrule
\end{tabular}
\end{center}

% ===================== 10 =====================
\section{Die Experimente und ihre Grenzen}
\label{sec:experimente}

\textbf{Das Bild.} Zwei Gläser stehen auf einem Tisch --- Gebiet A und Gebiet B.
Daneben das Weltmeer. Gebietsreduktion heißt: Eines der Gläser wird ins Weltmeer
geschüttet. Das Glasvolumen verschwindet nicht --- es verteilt sich auf 1,4
Milliarden Kubikkilometer. Liouville ist erfüllt: Kein Tropfen geht verloren.

Aber: Welches Glas es war --- A oder B --- ist im Rauschen des Ozeans so fein
verteilt, dass kein reales Instrument es je wieder herauslösen kann. Nicht weil
die Information vernichtet wäre (sie steckt in jeder Strömung, jedem
Temperaturgradienten), sondern weil die Auflösungsgrenze jedes realen Beobachters
zehn Größenordnungen unter dem Signal liegt.

Und: Das Meer kann danach unter denselben äußeren Bedingungen nicht dieselbe
Darstellungskapazität für die zwei Gläser bieten. Es kann sie nicht mehr separat
zeigen --- das Zielgebiet Z ist kleiner als $A+B$. Das ist nicht Unwissenheit; das
ist eine physikalische Tatsache über Gebietsgrößen.

Der Preis des Einschüttens: $Q \ge k_B T \ln 2$. Nicht wegen Unwissenheit ---
weil das Meer das Volumen aufnehmen musste und dafür Wärme floss.

\subsection{Die Größenordnungen}

\LK \pruef{Checks 2--3} Ein realer CMOS-Schaltvorgang liegt bei
${\sim}10^{-17}$--$10^{-16}$\,J --- vier bis fünf Größenordnungen über Landauer.
Alle Bit-Löschungen der Welt zusammen hätten Landauer-Kosten im Watt- bis
Kilowatt-Bereich (je nach Ansatz der globalen Löschrate, im Skript deklariert);
tatsächlich verbrauchen Rechenzentren einige zehn Gigawatt [B8] --- Faktor
$10^{7}$--$10^{10}$. In jeder realen Maschine ist die Grenze unsichtbar; sie ist
Naturgesetz-Untergrenze, der Rest ist Ingenieursverlust
($CV^2$-Umladung, Leckströme, Datentransport).

\LK \textbf{Das Weltmeer-Bild.} Das Bit-Volumen, das bei der Gebietsreduktion ins
Bad kippt, steht in einem so extremen Missverhältnis zum Bad (1 Freiheitsgrad
gegen $10^{23}$), dass es \textbf{unterhalb der Auflösungsgrenze} jedes realen
Instruments liegt. Das Bad nimmt das Volumen auf --- Liouville ist erfüllt ---
aber kein Messgerät kann das Signal noch vom Rauschen des Bades trennen. Nicht
weil das Instrument unvollkommen wäre, sondern weil die Struktur des Systems
selbst die Auflösungsgrenze setzt.

\subsection{Epistemisch vs.\ ontologisch --- ein notwendiger Vorbehalt}

\LB Nicht messbar heißt nicht nicht vorhanden. Die Mikrozustände des Bades nach
der Gebietsreduktion existieren alle noch; jedes Teilchen hat einen genauen Ort
und Impuls; die Information, welches Ausgangsgebiet (A oder B) das System hatte,
steckt noch in den Korrelationen --- Liouville garantiert das ausdrücklich. Das
Bad hat das Volumen aufgenommen, nicht vernichtet.

Landauer gilt \textbf{nicht, weil} Information ontologisch vernichtet wird. Es
gilt, \textbf{weil} das Bad das Volumen aufnehmen muss (Liouville) --- das ist
eine Aussage über die Mechanik, nicht über Wissen oder Wahrnehmung. Ein realer
Beobachter kann die im Bad verteilten Korrelationen nicht auflösen, aber das ist
die epistemische Seite; die thermodynamische Seite (Wärmefluss) folgt aus
Liouville allein.

\LB \textbf{Was aber ontologisch wahr ist:} Ein kleineres Gebiet kann unter
denselben Bedingungen weniger Zustände tragen als ein größeres. Nach der
Gebietsreduktion $\{A,B\} \to Z$ gilt $W_Z \le W_{A+B}/2$ --- das Zielgebiet ist
höchstens halb so groß wie das Ausgangsgebiet. Es kann daher unter denselben
äußeren Bedingungen (selbe Temperatur, selber Druck, selbe Kopplung ans Bad) nicht
dieselbe Menge an unterscheidbaren Zuständen tragen wie vorher. Das ist keine
Unwissenheit des Beobachters, sondern eine physikalische Tatsache über
Gebietsgrößen. Wer nach der Gebietsreduktion unter denselben Bedingungen dasselbe
darstellen will wie vorher, braucht ein anderes --- größeres --- physikalisches
System.

Die Korrelationen, die im Bad verteilt sind, sind real. Sie sind nur so fein
verteilt, dass kein Beobachter mit endlichen Mitteln und unter denselben
Bedingungen sie rekonstruieren kann. Das ist der Übergang vom ontologischen zum
epistemischen: Das Gebiet ist kleiner (ontologisch), und seine Rekonstruktion
würde ein Bad-großes Wissen erfordern (epistemisch).

\LK \textbf{Maxwells Dämon [B12] und Bennetts Auflösung [B3].} Ein Wesen, das alle
$10^{23}$ Mikrozustände des Bades kennt --- ein Laplace-Dämon ---, könnte das Bit
in Gedanken zurückverfolgen. Für ihn wäre Löschen reversibel und würde nichts
kosten. Maxwells Dämon zeigt genau das: Ein Beobachter mit vollständigem
Mikrozustandswissen kann Entropie scheinbar ohne Kosten verringern. Bennetts
Auflösung 1982 [B3]: Der Dämon muss seine \textbf{Messnotizen löschen}, um einen
neuen Zyklus zu beginnen --- und genau das kostet $k_B T \ln 2$. Der eigentliche
Ort der Entropieproduktion ist das Löschen der Notizen, nicht die Messung selbst.
Damit bleibt der zweite Hauptsatz erhalten, aber die Begründung ist präziser:
nicht „Information geht verloren“, sondern „Gebietsreduktion von Information
erzeugt Entropie“.

Zusammengefasst: Landauer ist \textbf{ontologisch} durch Liouville begründet. Die
\textbf{Irreversibilität} ist statistisch, nicht ontologisch absolut --- für einen
Dämon mit vollständigem Wissen wäre sie aufhebbar. Für jeden realen Beobachter ist
die Auflösungsgrenze jedoch so tief unter dem Signal, dass die statistische
Irreversibilität praktisch absolut ist. Beides muss gesagt werden.

\subsection{Die Einzelteilchen-Experimente}

Bérut et al.\ 2012 [B4] (Kolloid in optischer Doppelmulde, Zykluszeiten
10--40\,s), Jun/Gavrilov/Bechhoefer 2014 [B5] (Feedback-Falle, höhere Präzision).
Beide nähern sich der Grenze und bestätigen die $1/\tau$-Form der
Zusatzdissipation \pruef{Check 9}. Sie arbeiten im quasistatischen Grenzfall ---
kleinstes mögliches Bad, extrem langsame Prozesse ---, um überhaupt in die Nähe
der Grenze zu kommen. Ein realer Rechner löscht Bits in einer Nanosekunde ins
Weltmeer, im Sturm: Die Verwirbelung dominiert alles, die Grenze liegt zehn
Größenordnungen tiefer als der Lärm.

\textbf{Drei deklarierte Skepsis-Punkte} (Einwände gegen die Beweiskraft der
Experimente, nicht gegen das Prinzip --- das Prinzip hängt an
Abschnitt~\ref{sec:liouville}):

\begin{enumerate}[leftmargin=1.6em,itemsep=0.25em]
  \item \textbf{Extrapolation:} Bestätigt wird ein Limes ($\tau\to\infty$), den
        kein Experiment betritt; die Grenze wird als Asymptote gefittet.
  \item \textbf{Signal $\approx$ Rauschen:} $k_B T \ln 2$ liegt in der
        Größenordnung der thermischen Fluktuation einer Einzelmessung; das
        Resultat existiert nur als Ensemble-Mittel; manche Einzelzyklen
        „gewinnen“ Energie.
  \item \textbf{Erfolgsquote:} Die Protokolle löschen mit ${\sim}90$--95\,\%
        Zuverlässigkeit; die Verbuchung fehlgeschlagener Löschungen geht in die
        Bilanz ein.
\end{enumerate}

\LK \textbf{Feedback-Erweiterung.} Mit Messung und Rückkopplung gilt der
verallgemeinerte zweite Hauptsatz $W_{\text{ext}} \le -\Delta F + k_B T\,I$
(Sagawa--Ueda [B9]), experimentell bestätigt [B10][B11]. Information ist in der
Buchhaltung eine Ressource mit Energie-Wechselkurs $k_B T$ pro nat
\pruef{Check 10}.

% ===================== 11 =====================
\section{Anschluss an offene Fragen (intern)}

\begin{itemize}[leftmargin=1.4em,itemsep=0.35em]
  \item \textbf{Feedback-Buchhaltung in schnellen Systemen:} Rechenoperationen bei
        ${\sim}$ns-Zeitskalen liegen zehn Größenordnungen vom getesteten
        quasistatischen Regime entfernt; dort dominiert der $B/\tau$-Term alles.
        Wer in schnellen rekurrenten Systemen ein Landauer-artiges Residual messen
        will, muss den Sagawa-Ueda-Term $k_B T\,I$ [B9] explizit ins konventionelle
        Budget aufnehmen, bevor ein Residual als neu gelten kann --- sonst misst
        man Feedback-Thermodynamik, nicht neue Physik.
  \item \textbf{FFGFT-Berührung: keine.} Die Landauer-Buchhaltung lebt auf der
        Ebene der statistischen Mechanik; das $\xi$-Schema wird nicht berührt
        (eigene Ebene, eigene Buchhaltung --- P20/P39-analog).
\end{itemize}

% ===================== 12 =====================
\section{Kurzfassung in fünf Sätzen}

Gebietsreduktion verkleinert das Phasenraumgebiet des Systems; Liouville verbietet
Volumenschrumpfung in abgeschlossenen Systemen; also muss ein Bad das Volumen
übernehmen, und der Preis ist Wärme
$\ge k_B T \ln(W_{\text{v}}/W_{\text{n}})$ --- die $\ln 2$ nur als
Binär-Spezialfall. Gezählt wird nach Unterscheidbarkeit (Orthogonalität), nicht
nach Energie; die absolute Grund-Diskretisierung liefert $h^f$, alles Gröbere ist
emergente Zählung auf gewählter Ebene. Nicht messbar heißt nicht nicht vorhanden:
Die Mikrozustände des Bades existieren alle noch, die Information über das
gelöschte Bit steckt in Korrelationen --- Landauer gilt wegen der Mechanik
(Liouville), nicht wegen Unwissenheit; die Irreversibilität ist statistisch, nicht
ontologisch absolut (Maxwells Dämon, Bennetts Auflösung). Die Buchhaltung zählt
Gebiete und liest keine Inhalte --- jede Gebietsbelegung kostet identisch. Real
wird die Grenze nie erreicht, weil jede endliche Geschwindigkeit
Zusatzdissipation erzeugt, die um Größenordnungen dominiert --- das Bit kippt ins
Weltmeer im Sturm, die Auflösungsgrenze des Bades liegt zehn Größenordnungen über
der Landauer-Grenze.

% ===================== 13 =====================
\section{Schichtstatus-Zusammenfassung}

Die folgende Tabelle ist eine \emph{Auswahl} der tragenden Aussagen; im Fließtext
sind weitere Stellen markiert.

\begin{center}
\small
\begin{longtable}{@{}p{0.72\textwidth}l@{}}
\toprule
\textbf{Aussage} & \textbf{Status} \\
\midrule
\endfirsthead
\toprule
\textbf{Aussage} & \textbf{Status} \\
\midrule
\endhead
Bit $\ne$ physikalisches Gebiet (Shannon vs.\ Boltzmann) & \LB \\
$S = k_B \ln W$ als Ausgangspunkt & \LS \\
Gebietsreduktion $\{A,B\}\to Z$ (Physik, keine Beschreibung) & \LS \\
Gebietsreduktion $\Rightarrow \Delta S = -k_B \ln(W_{A+B}/W_Z)$ & \LB \\
Liouville-Theorem & \LB \\
Hinreichend ist die Nicht-Bijektivität der Gebietsabbildung, nicht die logische Irreversibilität & \LB \\
Abgeschlossenes System kann nicht löschen & \LB \\
Landauer-Grenze $Q \ge k_B T \ln 2$ & \LB \\
Numerischer Wert $2{,}87\cdot10^{-21}$\,J bei 300\,K & \LK \\
Zählbasis $=$ Orthogonalität, nicht Energie & \LB \\
$h^f$ als Grund-Diskretisierung & \LB \\
Gibbs-Form als allgemeine Entropie & \LS \\
Inhalt (wahr/falsch) geht nicht in die Rechnung ein & \LB \\
Hardware $\ne$ Rechenoperation & \LB \\
Rechner schrumpft nicht & \LB \\
Korrelation ins Bad (nicht Silizium) & \LB \\
Vergessensvorgang inhaltsneutral & \LB \\
DRAM ${\sim}1{,}12\cdot10^{5}$ Refresh-Zyklen pro Stunde (32\,ms) & \LK \\
No-Cloning-Theorem & \LB \\
$T_2 \sim 100\,\text{\textmu s}$ supraleitende Qubits & \LK \\
Minimalpreis Kohärenzerhaltung & \LH \\
Volumenerhaltende Abbildungen: keine Untergrenze & \LB \\
Volumenverkleinernde Abbildungen: das Bad zahlt & \LB \\
Epistemisch vs.\ ontologisch: Korrelationen bleiben & \LB \\
Sagawa-Ueda-Term $k_B T\,I$ & \LK \\
\bottomrule
\end{longtable}
\end{center}

\textbf{Bilanz der Tabelle:} 17$\times$ \LB (algebraisch/mathematisch aus
deklarierten Grundlagen) $\cdot$ 4$\times$ \LK (empirisch geprüft bzw.\
Primärquellen) $\cdot$ 3$\times$ \LS (deklarierte Ausgangspunkte) $\cdot$
1$\times$ \LH (offene Forschungsfrage) $=$ 25 Einträge.

% ===================== ANHANG KORREKTUREN =====================
\clearpage
\phantomsection
\addcontentsline{toc}{section}{Anhang: Korrekturen gegenüber der Arbeitsnotiz}
\section*{Anhang: Korrekturen gegenüber der Arbeitsnotiz}
\label{sec:korr}

Diese A-Serien-Fassung folgt der Arbeitsnotiz vom 24.~Juli 2026 inhaltlich
vollständig. Die folgenden Punkte wurden dabei geändert und sind hier offen
ausgewiesen, damit die Änderungen prüfbar bleiben.

\begin{enumerate}[leftmargin=1.6em,itemsep=0.4em]
  \item \textbf{Zustandszahl im Kolloid-Gebiet vereinheitlicht.} Die Notiz nannte
        an einer Stelle $7{,}9\cdot10^{9}$ und an anderer ${\approx}10^{5}$
        orthogonale Zustände, beide unter Berufung auf Check~7. Hier steht
        durchgehend $7{,}9\cdot10^{9}$. \emph{Zu prüfen:} welcher Wert im
        Prüfskript tatsächlich steht.
  \item \textbf{Dekohärenzzeit vs.\ Refresh-Intervall.} Die Notiz gab
        $T_2 \sim 100\,\text{\textmu s}$ als „siebenmal kürzer“ als 64\,ms an;
        korrekt sind rund 640-mal. Korrigiert.
  \item \textbf{DRAM-Raten getrennt.} Die Notiz mischte ${\sim}15{.}000$
        Refresh-Vorgänge pro Sekunde mit ${\sim}112{.}000$ Zyklen pro Stunde und
        Zelle; beides zusammen ist um Faktor $10^{3}$ inkonsistent. Hier
        getrennt: pro Zelle rund $1{,}12\cdot10^{5}$ Zyklen pro Stunde bei
        32\,ms Intervall (bei 64\,ms rund $5{,}6\cdot10^{4}$); der Controller
        gibt bei 8192 Zeilen und 64\,ms rund $1{,}3\cdot10^{5}$ Refresh-Befehle
        pro Sekunde und Bank aus.
  \item \textbf{Schichtbilanz korrigiert.} Die Notiz zählte
        15$\times$\,[B] und 8$\times$\,[K]; die Tabelle enthält tatsächlich
        16$\times$\,[B] und 4$\times$\,[K] bei 24 Einträgen (mit dem
        Zusatzeintrag aus Punkt~11 jetzt 17$\times$\,[B] bei 25 Einträgen). Ausgezählt und
        korrigiert; zugleich als \emph{Auswahl} deklariert, da im Fließtext
        weitere Stellen markiert sind.
  \item \textbf{Quellenliste ergänzt.} [B11], [B12], [B13] und [B20] waren im Text
        zitiert, aber nicht aufgeführt; der Eintrag [B17] enthielt versehentlich
        die Texte von [B13] (Fredkin/Toffoli) und [B12] (Maxwells Dämon). Alle
        vier ergänzt, [B17] bereinigt. Die rekonstruierten Einträge sind in der
        Quellenliste als solche markiert und vor der Veröffentlichung gegen die
        Arbeitsunterlagen zu prüfen.
  \item \textbf{Marker [H].} Das A-Serien-Schema kennt [SETZUNG], [K], [B] und
        [S]. Die Notiz verwendet zusätzlich [H] für „offene Forschungsfrage“.
        Hier beibehalten und in der Legende deklariert --- \emph{offen}, ob [H]
        ins Schema aufgenommen oder die Stelle nach [S] überführt wird.
  \item \textbf{Reihenfolge.} Die Notiz führte die Abschnitte 6c vor 6b. Hier in
        die logische Reihenfolge gebracht (Hardware $\ne$ Rechenoperation, dann
        Hardware-Seite, dann Dauerbetrieb, dann Quantencomputer) und als
        Unterabschnitte von Abschnitt~\ref{sec:inhalt} geführt.
  \item \textbf{Abbildungen.} Die Kasten- und Pfeildiagramme der Notiz verwendeten
        Unicode-Zeichnungszeichen; hier in reines ASCII gesetzt, damit der Bau
        glyphfrei bleibt. Inhalt unverändert.
  \item \textbf{Fehlender Absatzkopf.} Der Absatz zu den Größenordnungen
        (CMOS, Rechenzentren) begann in der Notiz ohne Kopf; hier als
        Unterabschnitt „Die Größenordnungen“ mit \LK gefasst.
  \item \textbf{Sprachliches.} Einzelne Grammatikfehler stillschweigend
        korrigiert (u.\,a.\ „das Informationsgehalt“, „Magnetfeldfluktuiation“,
        „was kostet das Gebietsreduktion“).
  \item \textbf{„Operation“ durch „Gebietsabbildung“ ersetzt (Landauer-Absatz).}
        Die Notiz schrieb: die logische Irreversibilität sei „eine hinreichende
        Bedingung für Wärmeproduktion bei einer bestimmten Klasse von
        \emph{Operationen}“. Das widerspricht dem eigenen Programm des Dokuments
        --- eine Operation ist eine Beschreibung, und Beschreibungen kosten
        nichts. Hinreichend ist allein die Nicht-Bijektivität der
        \emph{Gebietsabbildung}. Zugleich ist die Aussage schärfer gefasst: mit
        der Gebietsabbildung fällt die logische Irreversibilität nur unter einer
        deklarierten Realisierungsannahme zusammen (disjunkte Gebiete, kein
        mitgeführtes Gebiet); ohne sie trägt sie nichts, weil Bennett jede logisch
        irreversible Funktion bijektiv ausführt. Der Absatz war zudem
        durchgehend \LK markiert, obwohl nur die Landauer-Zitatstelle eine
        Primärquellen-Aussage ist; die abgeleitete Hälfte trägt jetzt \LB. Ein
        entsprechender Eintrag ist in die Schichtstatus-Tabelle aufgenommen.
\end{enumerate}

\textbf{Nicht geändert:} sämtliche physikalischen Aussagen, die Schichtzuordnungen
der einzelnen Stellen, die Formeln, die Argumentationsführung und der Befund
„FFGFT-Berührung: keine“.

% ===================== BELEGE =====================
\clearpage
\phantomsection
\addcontentsline{toc}{section}{Belege}
\section*{Belege}

\begingroup
\small
\begin{description}[leftmargin=3.2em,style=sameline,itemsep=0.3em]

\item[{[B1]}] R. Landauer, \emph{Irreversibility and Heat Generation in the
Computing Process}, IBM J.\ Res.\ Dev.\ \textbf{5}, 183 (1961).

\item[{[B2]}] Liouville-Theorem: Standardresultat der Hamiltonschen Mechanik;
z.\,B.\ Landau/Lifschitz, \emph{Mechanik}, §46; quantenmechanisches Gegenstück:
Unitarität der Zeitentwicklung.

\item[{[B3]}] C.~H. Bennett, \emph{The Thermodynamics of Computation --- a
Review}, Int.\ J.\ Theor.\ Phys.\ \textbf{21}, 905 (1982); ders., \emph{Notes on
Landauer's principle}, Stud.\ Hist.\ Phil.\ Mod.\ Phys.\ \textbf{34}, 501 (2003).

\item[{[B4]}] A. Bérut et al., \emph{Experimental verification of Landauer's
principle}, Nature \textbf{483}, 187 (2012).

\item[{[B5]}] Y. Jun, M. Gavrilov, J. Bechhoefer, \emph{High-Precision Test of
Landauer's Principle in a Feedback Trap}, Phys.\ Rev.\ Lett.\ \textbf{113},
190601 (2014).

\item[{[B6]}] Gibbs-/Shannon-Entropie: Standard; z.\,B.\ E.~T. Jaynes,
\emph{Information Theory and Statistical Mechanics}, Phys.\ Rev.\ \textbf{106},
620 (1957).

\item[{[B7]}] Semiklassische Zustandszählung $\Gamma/h^f$: Standard; z.\,B.\
Landau/Lifschitz, \emph{Statistische Physik}, §7.

\item[{[B8]}] Rechenzentrums-Energieverbrauch: IEA-Berichte 2024/25,
Größenordnung mehrere hundert TWh/Jahr weltweit ($\approx$ einige zehn GW
Dauerleistung). Zahl dient nur dem Größenordnungsvergleich.

\item[{[B9]}] T. Sagawa, M. Ueda, \emph{Second Law of Thermodynamics with
Discrete Quantum Feedback Control}, Phys.\ Rev.\ Lett.\ \textbf{100}, 080403
(2008); dies., \emph{Generalized Jarzynski Equality under Nonequilibrium Feedback
Control}, Phys.\ Rev.\ Lett.\ \textbf{104}, 090602 (2010).

\item[{[B10]}] S. Toyabe et al., \emph{Experimental demonstration of
information-to-energy conversion}, Nature Physics \textbf{6}, 988 (2010).

\item[{[B11]}] \emph{(rekonstruiert)} J.~V. Koski et al., \emph{Experimental
observation of the role of mutual information in the nonequilibrium dynamics of a
Maxwell demon}, Phys.\ Rev.\ Lett.\ \textbf{113}, 030601 (2014).

\item[{[B12]}] \emph{(rekonstruiert)} Maxwells Dämon --- Ursprung des
Gedankenexperiments: J.~C. Maxwell, \emph{Theory of Heat}, Longmans (1871);
Standarddarstellung in H.~S. Leff und A.~F. Rex (Hrsg.), \emph{Maxwell's Demon 2},
IOP Publishing (2003).

\item[{[B13]}] \emph{(rekonstruiert)} E. Fredkin und T. Toffoli, \emph{Conservative
Logic}, Int.\ J.\ Theor.\ Phys.\ \textbf{21}, 219 (1982) --- reversible Gatter und
die Frage des kontinuierlichen Entropiestroms physikalischer Rechner.

\item[{[B14]}] W.~K. Wootters und W.~H. Zurek, \emph{A single quantum cannot be
cloned}, Nature \textbf{299}, 802 (1982) --- No-Cloning-Theorem.

\item[{[B15]}] P.~W. Shor, \emph{Scheme for reducing decoherence in quantum
computer memory}, Phys.\ Rev.\ A \textbf{52}, R2493 (1995); A.~M. Steane,
\emph{Error correcting codes in quantum theory}, Phys.\ Rev.\ Lett.\ \textbf{77},
793 (1996) --- Quantenfehlerkorrektur.

\item[{[B16]}] F. Arute et al.\ (Google AI Quantum), \emph{Quantum supremacy using
a programmable superconducting processor}, Nature \textbf{574}, 505 (2019) ---
typische $T_2$-Zeiten supraleitender Qubits ${\sim}100\,\text{\textmu s}$.

\item[{[B17]}] M. Esposito und C. Van den Broeck, \emph{Three faces of the second
law: I.\ Master equation formulation}, Phys.\ Rev.\ E \textbf{82}, 011143 (2010);
P. Kammerlander und J. Anders, \emph{Coherence and measurement in quantum
thermodynamics}, Sci.\ Rep.\ \textbf{6}, 22174 (2016) --- Quantenthermodynamik und
Minimalpreis der Kohärenzerhaltung.

\item[{[B18]}] C.~E. Shannon, \emph{A Mathematical Theory of Communication}, Bell
System Technical Journal \textbf{27}, 379 (1948) --- Definition des Bit als
$\log_2$ der unterscheidbaren Zustände; explizit informationstheoretisch, nicht
thermodynamisch.

\item[{[B19]}] E.~T. Jaynes, \emph{Information Theory and Statistical Mechanics},
Phys.\ Rev.\ \textbf{106}, 620 (1957) --- Verbindung Shannon-Entropie zu
Boltzmann: gilt exakt nur bei Gleichverteilung über physikalische Mikrozustände;
bei grob gewählter logischer Auflösung nicht.

\item[{[B20]}] \emph{(rekonstruiert)} Johnson-Nyquist-Rauschen: J.~B. Johnson,
\emph{Thermal Agitation of Electricity in Conductors}, Phys.\ Rev.\ \textbf{32},
97 (1928); H. Nyquist, \emph{Thermal Agitation of Electric Charge in Conductors},
Phys.\ Rev.\ \textbf{32}, 110 (1928).

\end{description}
\endgroup

\vspace{1em}
\noindent
{\small\itshape Die als \emph{(rekonstruiert)} markierten Einträge fehlten in der
Arbeitsnotiz und wurden aus dem Zitationskontext ergänzt; sie sind vor der
Veröffentlichung gegen die Arbeitsunterlagen zu prüfen.}

\end{document}
